\SourceLeaf{62}
Der erste Teil der gestellten Aufgabe ist erledigt, da die Koeffizienten von $(1+n\cos\varphi)^{-1}$ durch eine endliche Formel bestimmt sind. Im verbleibenden zweiten Teil ist zu untersuchen, wie man von den gegebenen Koeffizienten einer beliebigen Potenz zu den gesuchten Koeffizienten der nächstniedrigeren Potenz gelangt.

$\delta$) Am einfachsten nimmt man die Ausdrücke für den Anfangskoeffizienten und den allgemeinen Koeffizienten der unbestimmten Reihe für $(1+n\cos\varphi)^{\nu}$, ersetzt $\nu$ durch $\nu-1$ und leitet dann die Beziehung der so erhaltenen Koeffizienten zu den vorherigen her.

I. Nach (6) war das Anfangsglied des in eine Reihe entwickelten Anfangskoeffizienten $A$ gleich $1$ und wird also von der Ersetzung nicht berührt. Das allgemeine Glied mit Index $h^{\text{tus}}$ war
\[
= \frac{1}{2^{2h}}\,{}^{\nu}\mathfrak{B}^{(2h)} .  {}^{2h}\mathfrak{B}^{(h)}\,n^{2h},
\]
und geht durch Einsetzen von $\nu-1$ anstelle von $\nu$ über in
\[
\frac{1}{2^{2h}}\,{}^{\nu-1}\mathfrak{B}^{(2h)} .  {}^{2h}\mathfrak{B}^{(h)}\,n^{2h}.
\]
Der zweite Ausdruck entsteht also offenbar aus dem ersten durch Multiplikation mit
\[
\frac{\nu-2h}{\nu} = 1-\frac{2h}{\nu}
\]
Das allgemeine Glied von $A$ mit Index $h^{\text{tum}}$ mit $\frac{2h}{\nu}$ zu multiplizieren ist gleichbedeutend damit, es nach $n$ zu differenzieren, durch $\nu\,\mathrm{d}n$ zu dividieren und mit $n$ zu multiplizieren. Bezeichnen wir also mit $B$ das Anfangsglied von $(1+n\cos\varphi)^{\nu-1}$, so gilt allgemein
\[
B = A - \frac{n\,{}^{*}\!\mathrm{d}A}{\nu\,\mathrm{d}n}.
\]
\par\noindent\hspace*{\parindent}II. Für $(1+n\cos\varphi)^{\nu}$ lautete das Anfangsglied des unbestimmten Koeffizienten $A^{(r)}$
\[
\frac{1}{2^{r-1}}\,{}^{\nu}\mathfrak{B}^{(r)}\,n^{r},
\]
Bezeichnen wir mit $B^{(r)}$ den zu $(1+n\cos\varphi)^{\nu-1}$ gehörigen Koeffizienten, der aus dem vorherigen durch Einsetzen von $\nu-1$ anstelle von $\nu$ entsteht, so lautet sein Anfangsglied
\[
\frac{1}{2^{r-1}}\,{}^{\nu-1}\mathfrak{B}^{(r)}\,n^{r}.
\]
Dieses Glied entsteht nun durch Multiplikation des vorherigen mit $\frac{\nu-r}{\nu}$ oder, was dasselbe ist,
